\begin{corollary}[unrestricted toroidal topology]\hypertarget{tag-02ZV}{}\label{push:src:A.4.5}\label{push:A45}
\StatementRecord{push:src:A.4.5}{proved}{02ZV}{\SeriesRef{push:src:A.2.4}, \SeriesRef{push:src:A.4.4}, \SeriesRef{push:src:A.4.7}, \SeriesRef{push:src:A.4.8}, \SeriesRef{push:src:A.4.9}}

For every output of \ref{push:src:A.2.4}, \ref{push:src:A.4.7}--\ref{push:src:A.4.9} determine the
Euler number, fundamental group, all integral homology
groups with torsion, nearby Euler function and marked
monodromy. The integral normalization computation is
given by the finite algorithm and the primitive
central pushforward matrices of \ref{push:src:A.4.9}; the independent
ambient computation is \ref{push:src:A.4.8}, and the torus-stratum
Euler count is \ref{push:src:A.4.7}. The non-K\"ahler limitation on
classical Clemens--Schmid is the explicit one in
\ref{push:src:A.4.4}; no polarized limiting Hodge structure is
asserted.
\end{corollary}

\begin{proof}
\begin{enumerate}[label=\arabic*.,leftmargin=*]
\item An output of \ref{push:src:A.2.4} has finitely many ray and
cone orbits with specified integral maps. Count its pure top
rays $k$, mixed rays $m$ and maximal four-cones $t_4$.
Theorem~\ref{push:src:A.4.7} gives $\chi=11k+t_4$ and
$\pi_1=\Z$, and evaluates the nearby Euler function on
each branch stratum; its integral is two.

\item Theorem~\ref{push:src:A.4.8} gives the full free integral
homology table and the marked degree-three monodromy $I+nN$,
with six nonzero Smith factors $n,n,n,n,n,11n$. The integral
splitting of \ref{push:src:A.4.9} computes the same groups
directly along the finite sequence of normalization modifications;
its exceptional-fibre diagonal entries are units, so it also
determines primitivity and torsion. Thus every output named in
the statement is determined by the indicated finite data.
The exact failure of the particular polarized comparison is
already the calculation in \ref{push:src:A.4.4}; no additional
Hodge-theoretic conclusion is used here.
\end{enumerate}
\end{proof}